\documentclass[11pt]{article}
\usepackage[utf8]{inputenc}
\usepackage[T1]{fontenc}
\usepackage{amsmath, amssymb, amsthm}
\usepackage{hyperref}
\usepackage{graphicx}
\usepackage{listings}
\usepackage{enumitem}
\usepackage{geometry}

\geometry{margin=1in}

\title{Quantum API Mermin Pérez\newline Documentation Bundle}
\author{Repository documentation}
\date{\today}

\begin{document}
\maketitle

\section{Overview}
This document captures the Markdown and LaTeX-oriented documentation bundle for the
Julia workspace. The source content is aligned with the local Documenter setup and
focuses on matrix-valued distributions and the solver workflow.

\section{Matrix-variate Wishart workflow}
The matrix-variate Wishart workflow follows the model
\[
X \sim \mathcal{W}_p(n, \Sigma)
\]
where $p = 8$ and $\Sigma$ is the positive-definite scale matrix.

\begin{itemize}[leftmargin=*]
    \item Build a scale matrix $\Sigma$.
    \item Construct the distribution as \texttt{Wishart(9, $\Sigma$)}.
    \item Draw samples with \texttt{rand}.
    \item Normalize the sample by its trace.
    \item Evaluate the score with \texttt{logpdf}.
\end{itemize}

\section{Distributions references}
The project documentation is anchored to the JuliaStats references:
\begin{itemize}[leftmargin=*]
    \item \url{https://juliastats.org/Distributions.jl/stable/}
    \item \url{https://juliastats.org/Distributions.jl/stable/matrix/}
    \item \url{https://juliastats.org/Distributions.jl/stable/mixture/}
    \item \url{https://juliastats.org/Distributions.jl/stable/fit/}
\end{itemize}

\section{Solver reference}
The repository also includes a 25×25 Sudoku solver pattern that can be documented
and rendered as part of the narrative site. The constraint system is based on the
usual Latin-square rules:
\[
\forall r,c:\quad \text{cell}(r,c) \in \{1, \dots, 25\}
\]

\section{Minimal Julia example}
\begin{lstlisting}[basicstyle=\ttfamily\small]
using Distributions
using LinearAlgebra

M = Matrix{Float64}(I, 2, 2)
println(tr(M))
\end{lstlisting}

\end{document}
